% ============================================================
%  TD Chapitre 2 : Equations et inequations
% ============================================================
\clearpage
\chapter{Équations et inéquations}
\chaptermark{Équations et inéquations}

\begin{exercice}[équations du premier degré et équations produit]
Résoudre dans $\R$ les équations suivantes :
\begin{enumerate}
\item $5(x - 2) = 3x + 4$ ;
\item $(3x - 2)(x + 5) = 0$ ;
\item $x^{2} = 4x$, résolue en factorisant (et non en divisant par $x$).
\end{enumerate}
\end{exercice}

\begin{exercice}[équations du second degré : le discriminant]
Résoudre dans $\R$ les équations suivantes en calculant le discriminant, puis factoriser quand c'est possible :
\begin{enumerate}
\item $2x^{2} - 5x - 3 = 0$ ;
\item $4x^{2} - 12x + 9 = 0$ ;
\item $x^{2} + x + 1 = 0$.
\end{enumerate}
\end{exercice}

\begin{exercice}[somme, produit et formation d'une équation]
\begin{enumerate}
\item calculer, si possible, la somme $S$ et le produit $P$ des solutions de $5x^{2} - 6x + 1 = 0$, puis vérifier en résolvant complètement l'équation ;
\item former l'équation du second degré dont les solutions ont pour somme $S = 2$ et pour produit $P = -8$, puis la résoudre ;
\item l'équation $2x^{2} - 7x + 3 = 0$ admet $\dfrac{1}{2}$ comme solution : retrouver l'autre solution avec le produit, sans calculer le discriminant.
\end{enumerate}
\end{exercice}

\begin{exercice}[équations bicarrées]
Résoudre dans $\R$ les équations suivantes avec le changement de variable $X = x^{2}$ :
\begin{enumerate}
\item $x^{4} - 5x^{2} + 4 = 0$ ;
\item $x^{4} + x^{2} - 2 = 0$.
\end{enumerate}
\end{exercice}

\begin{exercice}[inéquations du premier et du second degré]
Résoudre dans $\R$ les inéquations suivantes, en donnant l'ensemble des solutions sous forme d'intervalle ou de réunion d'intervalles :
\begin{enumerate}
\item $3x - 7 \leqslant 0$ ;
\item $-4x + 10 > 0$ ;
\item $x^{2} - 4x - 5 \leqslant 0$ ;
\item $-2x^{2} + 8x - 6 > 0$.
\end{enumerate}
\end{exercice}
